%% Model benchmarks, robustness, and SOTA comparison (main text).

\subsection{GOES-16 Unsupervised Detection Benchmark}
\label{sec:goes16-validation}

GOES-16 ABI CMIPF was evaluated against INPE ground references for 2024-10-31 (76 INPE hotspots; 5{,}184 valid grid cells; 284 truth cells after $3\times3$ morphological dilation). Low F1 ($\leq 0.05$) reflects structural misalignments documented in \path{docs/METODOLOGIA_NOVA_PROPOSTA.md}: temporal window mismatch, thermal anomaly vs.\ active-fire semantics, 2~km GOES cells vs.\ point-like INPE records, and CMIPF brightness temperature vs.\ dedicated AF products.

Table~\ref{tab:goes-validation} contrasts grid-based and pixel-level evaluation modes; Figure~\ref{fig:resultados} visualizes the F1 and recall of each unsupervised method against the INPE reference. Combined persistence and consensus ensemble (Lines~B/E) yield zero detections but also zero false positives---suitable for high-stakes alert filtering when fused with FIRMS in the operational pipeline.

\figurainline[0.8\linewidth]{figures/resultados-experimentais.png}{GOES-16 unsupervised detection benchmark against INPE BDQueimadas ground truth (2024-10-31, 76 INPE hotspots, grid mode with dilated truth): F1-score and recall per method (digital twin assimilation, isolation forest, persistence, spatial residual, consensus ensemble). All methods score F1$\leq$0.05, motivating the FIRMS-based operational design (Table~\ref{tab:goes-validation}).}{fig:resultados}

\begin{resulttable}{GOES-16 vs.\ INPE validation (2024-10-31). Grid: dilated truth (284 cells); pixel: 15~km focal matching.}{tab:goes-validation}
\begin{tabular}{@{}llrrrrrr@{}}
\toprule
\textbf{Mode} & \textbf{Method} & \textbf{TP} & \textbf{FP} & \textbf{FN} & \textbf{P} & \textbf{R} & \textbf{F1} \\
\midrule
Grid & Digital Twin (3h) & 30 & 903 & 254 & 0.032 & 0.106 & 0.049 \\
Grid & Isolation Forest & 18 & 680 & 266 & 0.026 & 0.063 & 0.037 \\
Grid & Spatial Residual & 8 & 718 & 41 & 0.011 & 0.163 & 0.021 \\
Grid & Combined Persistence & 0 & 0 & 284 & 0.000 & 0.000 & 0.000 \\
Pixel & Cluster match & 4 & 152 & 72 & 0.026 & 0.053 & 0.034 \\
\bottomrule
\end{tabular}
\end{resulttable}

\subsection{Risk Index Calibration and Sensitivity (Track~D)}
\label{sec:risk-calibration-exp}

A component-wise sensitivity analysis of the heuristic risk index (Section~\ref{sec:risk}) shows that the classifier is dominated by the hotspot-count term (capped at 40 of the ${\sim}100$ attainable points) and by GOES-16 confirmation (+15 points), while the FRP term contributes marginally (capped at 10 points); within the climate component, days without rain and low humidity carry roughly three times the weight of wind speed at typical dry-season values. Severity classifications are therefore most sensitive to hotspot-count and drought-related inputs, which motivates the calibration below.

The heuristic risk classifier was calibrated against a 111{,}054-record INPE BDQueimadas snapshot (12/Jun/2026; 19{,}249 positive / 11{,}362 negative municipal-day samples, Jan/2024--May/2026). Labels are defined directly from the reference data: a municipal-day is \emph{positive} if it contains at least one INPE hotspot and \emph{negative} otherwise. Because the risk index takes the daily hotspot count as its dominant input (Eq.~1), a municipal-day can only exceed the severity threshold when hotspots are present, so precision equals 1.000 \emph{by construction} (false positives are structurally impossible) and must not be read as an independent validation result. The informative quantities here are recall and F1 --- how many fire days the index flags at the calibrated thresholds --- and the calibration accordingly maximizes F1. Table~\ref{tab:risk-calibration} reports baseline vs.\ optimized weights: optimized weights increase hotspot sensitivity (\texttt{w\_focos} 8$\rightarrow$12), drought penalty (\texttt{w\_clima\_seca} 1.5$\rightarrow$2.5), and hotspot cap (\texttt{max\_focos} 40$\rightarrow$50), with severity thresholds lowered by 5~points. Out-of-sample evaluation against labels independent of the index inputs (e.g., burned-area products such as MapBiomas Fogo) remains future work.

\begin{resulttable}{Risk index calibration on the INPE BDQueimadas snapshot (111{,}054 records, 12/Jun/2026). Precision is 1.000 by construction (see text); recall and F1 are the informative metrics.}{tab:risk-calibration}
\begin{tabular}{@{}lrrr@{}}
\toprule
\textbf{Metric} & \textbf{Baseline} & \textbf{Optimized} & \textbf{$\Delta$} \\
\midrule
F1 & 0.833 & \textbf{0.936} & +0.104 \\
Precision & 1.000 & 1.000 & --- \\
Recall & 0.713 & \textbf{0.880} & +0.167 \\
Accuracy & 0.820 & \textbf{0.925} & +0.105 \\
\bottomrule
\end{tabular}
\end{resulttable}

\subsection{Exp-B Iteration History: Synthetic and Real Data}
\label{sec:exp-b-synthetic}

\textbf{v1 (200 timesteps).} Initial Koopman-VAE and NeKo-PIGNN underperformed shallow baselines (R\textsuperscript{2}$\approx$0.28--0.34 vs.\ MLP 0.93) due to KL-regularized VAE competing with prediction (Table~\ref{tab:exp-v1}, Appendix). \textbf{v2 (500 timesteps).} Architecture fixes---deterministic Koopman, curriculum learning, spectral regularization---reversed the gap. Table~\ref{tab:regression_synthetic} shows NeKo-PIGNN v2 reaching R\textsuperscript{2}=0.972 on synthetic next-day prediction, establishing an upper bound before real-data noise.

\begin{resulttable}{Exp-B v2: next-day prediction on synthetic data (500 steps, 30 nodes). Best in \textbf{bold}.}{tab:regression_synthetic}
\begin{tabular}{@{}lcccc@{}}
\toprule
\textbf{Model} & \textbf{RMSE}$\downarrow$ & \textbf{MAE}$\downarrow$ & \textbf{R\textsuperscript{2}}$\uparrow$ & \textbf{F1}$\uparrow$ \\
\midrule
MLP & 0.069 & 0.024 & 0.968 & 0.975 \\
LSTM & 0.084 & 0.045 & 0.951 & 0.945 \\
XGBoost & 0.087 & 0.039 & 0.948 & 0.931 \\
Koopman-Det (ours) & 0.068 & 0.022 & 0.968 & 0.975 \\
\textbf{NeKo-PIGNN v2 (ours)} & \textbf{0.064} & 0.024 & \textbf{0.972} & 0.975 \\
\bottomrule
\end{tabular}
\end{resulttable}

\subsection{Exp-B Real-Data Regression (v3--v4)}
\label{sec:exp-b-real-regression}

\textbf{v3} validated all models on real FIRMS/INPE/Open-Meteo data. MLP led in RMSE/R\textsuperscript{2}; all models maintained binary F1$>$0.92 on next-day fire-count regression---a \emph{different task} from three-class alert metrics (Section~\ref{sec:three-level}). The synthetic-to-real R\textsuperscript{2} drop (Table~\ref{tab:real-regression-combined}b) ranges from 17 to 25 percentage points, expected under label noise and data scarcity with only 67 training days.

\textbf{v4} tested feature engineering with 62--67 training samples; MLP remained best (R\textsuperscript{2}=0.782), while feature expansion degraded MLP to R\textsuperscript{2}=0.20 (Table~\ref{tab:exp-v4}, Appendix). With ${\sim}$250K-parameter models and ${\sim}$62 samples, simpler architectures dominate---a finding consistent across the Exp-B iteration.

\begin{resulttable}{Real-data regression (v3) and synthetic-to-real R\textsuperscript{2} drop (v2$\rightarrow$v3).}{tab:real-regression-combined}
\begin{minipage}[t]{0.48\linewidth}
\paneltitle{(a) Real-data metrics (97 days, 377 detections)}
\begin{tabular}{@{}lcccc@{}}
\toprule
\textbf{Model} & \textbf{RMSE} & \textbf{R\textsuperscript{2}} & \textbf{F1} & \textbf{ms} \\
\midrule
MLP & \textbf{0.138} & \textbf{0.799} & 0.941 & 0.0 \\
LSTM & 0.163 & 0.718 & 0.943 & 0.4 \\
XGBoost & 0.149 & 0.766 & 0.931 & 5.6 \\
Koopman-Det & 0.157 & 0.739 & 0.924 & 0.1 \\
NeKo-PIGNN v2 & 0.152 & 0.756 & 0.926 & 0.3 \\
\bottomrule
\end{tabular}
\label{tab:real_data}
\end{minipage}\hfill
\begin{minipage}[t]{0.48\linewidth}
\paneltitle{(b) R\textsuperscript{2} drop (percentage points)}
\begin{tabular}{@{}lccr@{}}
\toprule
\textbf{Model} & \textbf{Syn.} & \textbf{Real} & \textbf{$\Delta$} \\
\midrule
MLP & 0.968 & 0.799 & $-$17 \\
LSTM & 0.951 & 0.718 & $-$25 \\
XGBoost & 0.948 & 0.766 & $-$19 \\
Koopman-Det & 0.968 & 0.739 & $-$24 \\
NeKo-PIGNN v2 & 0.972 & 0.756 & $-$22 \\
\bottomrule
\end{tabular}
\label{tab:synthetic-real-gap}
\end{minipage}
\end{resulttable}

\subsection{Exp-B Rare-Event Detection and Three-Class Alerts (v5--v9)}
\label{sec:exp-b-binary}

\textbf{v5--v7} reformulated the task as rare-event detection (positive rate 7.5\%), then added persistence priors (semi-arid fires persist 2--7 days). NeKo-PIGNN achieved the highest Average Precision (0.353) in v5; ensemble $\times$ persistence reached 47.5\% precision with 21~FP in v7 (\ref{app:exp-b}). \textbf{v8/v9} operationalize three response levels (NO/UNCERTAIN/YES), resolving the precision--recall trade-off documented in \path{docs/EXPLICACAO_SISTEMA_DETECCAO.md}.

\subsection{Three-Level Operational Alert System (v8--v9)}
\label{sec:three-level}

Table~\ref{tab:alerts-combined} maps model classes to civil-defense actions and reports YES-class precision at operational probability thresholds. Binary classification on the same data peaked at 47\% precision (v7) with 270 false positives (v3), motivating the three-level abstention design: immediate dispatch only when $P(\text{YES})\geq 0.30$, monitoring for UNCERTAIN, and satellite-only coverage otherwise.

\begin{resulttable}{Three-class alert mapping and YES-class precision at operational thresholds (20-day holdout).}{tab:alerts-combined}
\begin{minipage}[t]{0.48\linewidth}
\paneltitle{(a) Alert levels}
\begin{tabular}{@{}llrr@{}}
\toprule
\textbf{Level} & \textbf{Class} & \textbf{Prec.} & \textbf{Rec.} \\
\midrule
Alert & YES & 82.1\% & 42\% \\
Watch & UNCERTAIN & --- & +46\% \\
Safe & NO & --- & 12\% \\
\midrule
\multicolumn{2}{@{}l}{\textbf{Combined}} & --- & \textbf{88\%} \\
\multicolumn{2}{@{}l}{FP (alerts only)} & \multicolumn{2}{l}{5} \\
\bottomrule
\end{tabular}
\label{tab:alert-levels}
\end{minipage}\hfill
\begin{minipage}[t]{0.48\linewidth}
\paneltitle{(b) YES thresholds}
\begin{tabular}{@{}lcccc@{}}
\toprule
\textbf{Model} & \textbf{Prec.} & \textbf{Rec.} & \textbf{TP} & \textbf{FP} \\
\midrule
XGB, $P\geq 0.30$ & 82.1\% & 41.8\% & 23 & 5 \\
XGB, $P\geq 0.50$ & 79.2\% & 34.6\% & 19 & 5 \\
NeKo (YES) & 91.7\% & 12.7\% & 11 & 1 \\
Ensemble, $P\geq 0.60$ & 88.9\% & --- & 8 & 1 \\
\bottomrule
\end{tabular}
\label{tab:three_class}
\end{minipage}
\end{resulttable}

Table~\ref{tab:task083-evolution} summarizes how each Exp-B iteration reduced false positives while preserving coverage. Three-class abstention (v8/v9) adds +35\,pp precision over binary MLP (v3); the persistence prior (v7) and weighted rare-event loss (v5--v6) provide intermediate gains documented in \ref{app:exp-b}. Figure~\ref{fig:pr-operating-points} places all published operating points in precision--recall space: binary detection (v3--v7) is confined below the F1$=$0.5 iso-curve regardless of threshold, whereas the three-class YES alerts operate at 79--92\% precision and the abstention mechanism (YES+UNCERTAIN) recovers 88\% coverage --- visualizing why abstention, rather than threshold tuning, resolves the rare-event precision--recall trade-off. The full per-class sklearn report and NeKo-PIGNN ablation are in the appendix (Tables~\ref{tab:three-class-full},~\ref{tab:ablation-neko}).

\figurainline[0.82\linewidth]{figures/pr-operating-points.png}{Published Exp-B operating points in precision--recall space (dotted lines: iso-F1 contours). Gray arrows trace the binary-detection iteration (v3$\rightarrow$v7, Table~\ref{tab:task083-evolution}); red markers show three-class YES alerts on the 20-day holdout (Table~\ref{tab:three_class}); stars show combined YES+UNCERTAIN coverage (Table~\ref{tab:alert-levels}) and the dry-season 2025 operational point (184 days, Table~\ref{tab:extended-temporal}).}{fig:pr-operating-points}

\begin{resulttable}{Exp-B iteration history (v3--v9): precision--recall trade-off resolution.}{tab:task083-evolution}
\begin{tabular}{@{}llrrr@{}}
\toprule
\textbf{Ver.} & \textbf{Approach} & \textbf{Prec.} & \textbf{Rec.} & \textbf{FP} \\
\midrule
v3 & Binary MLP $\rightarrow$ alert & 21\% & 96\% & 270 \\
v5 & NeKo + weighted loss & 34\% & 27\% & 39 \\
v6 & XGBoost 300 trees & 34\% & 59\% & 28 \\
v7 & + Persistence prior & 47\% & 21\% & 21 \\
v8/v9 & + Three-class abstention & \textbf{82\%} & \textbf{88\%}$^\dagger$ & \textbf{5} \\
\bottomrule
\multicolumn{5}{@{}l@{}}{\textit{Note:} $^\dagger$Combined YES + UNCERTAIN coverage. NeKo YES-only: 91.7\% prec., 12.7\% rec., 1~FP.} \\
\end{tabular}
\end{resulttable}

\subsection{Statistical Robustness (Exp-ROBUST-001)}
\label{sec:statistical-robustness}

Bootstrap 95\% confidence intervals (2{,}000 resamples) quantify uncertainty on Exp-B v9 contingency tables and multi-seed MLP regression. XGBoost YES alert precision (82.1\%) has CI [67.9\%, 96.4\%] (23 TP, 5 FP); combined YES+UNCERTAIN coverage 88\% CI [84.6\%, 91.2\%]. Reproduced MLP regression yields RMSE $0.137 \pm 0.003$ and R\textsuperscript{2} $0.800 \pm 0.010$ [0.783, 0.811], consistent with Table~\ref{tab:real_data}. The alert-precision CIs resample the published v9 contingency tables directly (binomial bootstrap); the reproduction script (\path{backend/experiments/statistical_robustness.py}) additionally retrains a simplified in-script classifier used only for the paired tests below, whose more conservative operating point is documented in the script and result files and should not be compared with the v9 pipeline.

Paired significance tests complement the bootstrap analysis (Table~\ref{tab:robustness-combined}b). Wilcoxon signed-rank on per-seed MLP vs.\ XGBoost RMSE yields $p=0.0625$; McNemar tests for three-class and YES-alert vs.\ persistence are not significant at $\alpha=0.05$ on the 20-day holdout. We report these for transparency rather than as superiority claims.

\begin{resulttable}{Bootstrap 95\% CI and paired significance tests (Exp-ROBUST-001).}{tab:robustness-combined}
\begin{minipage}[t]{0.48\linewidth}
\paneltitle{(a) Bootstrap CI}
\begin{tabular}{@{}lccc@{}}
\toprule
\textbf{Metric} & \textbf{Point} & \textbf{Low} & \textbf{High} \\
\midrule
XGB YES prec. & 82.1\% & 67.9\% & 96.4\% \\
NeKo YES prec. & 91.7\% & 75.0\% & 100.0\% \\
Coverage & 88.0\% & 84.6\% & 91.2\% \\
MLP RMSE & 0.1372 & 0.1334 & 0.1429 \\
MLP R$^2$ & 0.800 & 0.783 & 0.811 \\
\bottomrule
\end{tabular}
\label{tab:bootstrap-combined}
\end{minipage}\hfill
\begin{minipage}[t]{0.48\linewidth}
\paneltitle{(b) Paired tests}
\begin{tabular}{@{}p{2.5cm}rrl@{}}
\toprule
\textbf{Test} & \textbf{$p$} & \textbf{Sig.} & \textbf{Notes} \\
\midrule
Wilcoxon RMSE & 0.0625 & No & $\Delta$=-0.0099 \\
McNemar 3-class & 0.1418 & No & acc 73.2\% \\
McNemar YES & 0.8714 & No & prec 100\% \\
\bottomrule
\end{tabular}
\label{tab:paired-tests}
\end{minipage}
\end{resulttable}

\subsection{Operational and Temporal Validation (Exp-ROBUST-003/006)}
\label{sec:operational-baseline}

Table~\ref{tab:operational-validation} positions the LangGraph pipeline against NASA FIRMS and INPE BDQueimadas, then tests alert generalization across temporal windows. Cold-start latency (52.4~s) is ${\sim}309\times$ faster than the 3--6~h FIRMS publication window. Dry-season three-class alerts reach 84.2\% precision and 91.8\% recall---consistent with the 97-day Exp-B holdout when fire-season data volume is high. Mixed 12-month windows yield 31.6\% precision, indicating season-specific calibration. Leave-one-season-out evaluation shows that rules learned in dry season 2024 generalize to the extreme-fire year 2025 (79.0\% precision, 85.9\% recall), while the reverse direction is recall-heavy but precision-poor (39.2\%), supporting season-aware deployment.

\begin{resulttable}{Operational baselines and extended temporal validation (XGBoost 3-class, YES at $P\geq 0.30$).}{tab:operational-validation}
\begin{minipage}[t]{0.48\linewidth}
\paneltitle{(a) System comparison}
\begin{tabular}{@{}p{1.8cm}p{1.2cm}cc@{}}
\toprule
\textbf{System} & \textbf{Lat.} & \textbf{Prec.} & \textbf{Rec.} \\
\midrule
NASA FIRMS & 3--6 h & 80.0\% & 15.7\% \\
INPE BD\-Queimadas & 3--6 h & --- & --- \\
LangGraph (ours) & 52 s & 84.2\% & 91.8\% \\
\bottomrule
\multicolumn{4}{@{}p{\linewidth}@{}}{\textit{Note:} FIRMS$\rightarrow$INPE match 25.8\% (1\,km, 8-day sample). Cost: \$20/mo.} \\
\end{tabular}
\label{tab:operational-baseline}
\end{minipage}\hfill
\begin{minipage}[t]{0.48\linewidth}
\paneltitle{(b) Temporal windows}
\begin{tabular}{@{}p{1.6cm}rrrr@{}}
\toprule
\textbf{Window} & \textbf{Days} & \textbf{Prec.} & \textbf{Rec.} & \textbf{FP} \\
\midrule
Short subset & 117 & 34.4\% & 38.9\% & 40 \\
Dry 2024 & 141 & 44.9\% & 66.4\% & 114 \\
Dry 2025 & 184 & 84.2\% & 91.8\% & 75 \\
Full 12 mo & 316 & 31.6\% & 74.1\% & 130 \\
LOO 24$\rightarrow$25 & 500 & 79.0\% & 85.9\% & 379 \\
LOO 25$\rightarrow$24 & 500 & 39.2\% & 98.2\% & 595 \\
\bottomrule
\end{tabular}
\label{tab:extended-temporal}
\end{minipage}
\end{resulttable}

\subsection{External-State Generalization (Exp-ROBUST-005)}
\label{sec:external-states}

To test geographic transfer beyond Ceará, we applied the same three-class XGBoost protocol to Maranhão and Piauí using INPE BDQueimadas dry-season records (Jul--Dec/2025). Ceará retains the best precision--recall balance (84.2\% / 91.8\%); neighboring states show higher fire volume and recall (${>}96$\%) but lower precision (${\sim}58$--60\%), consistent with class imbalance and semi-arid persistence patterns (Table~\ref{tab:external-states}).

\begin{resulttable}{External-state validation: XGBoost 3-class YES alert at $P\geq 0.30$ (dry season Jul--Dec/2025).}{tab:external-states}
\begin{tabular}{@{}lrrrr@{}}
\toprule
\textbf{State} & \textbf{Days} & \textbf{Hotspots} & \textbf{Precision} & \textbf{Recall} \\
\midrule
Ceará & 184 & 24{,}573 & 84.2\% & 91.8\% \\
Maranhão & 184 & 233{,}854 & 57.5\% & 99.5\% \\
Piauí & 183 & 145{,}220 & 60.1\% & 96.4\% \\
\bottomrule
\end{tabular}
\end{resulttable}

Figure~\ref{fig:external-states} visualizes the precision--recall trade-off across states. These results support a Northeast pilot deployment model: calibrate thresholds per state rather than transferring Ceará parameters directly to higher-volume neighbors.

\figurainline[0.85\linewidth]{figures/mapa-external-states.png}{Dry-season 2025 alert precision and recall: Ceará vs.\ Maranhão and Piauí (INPE BDQueimadas).}{fig:external-states}

\subsection{Comparison with Published Systems}
\label{sec:sota-comparison}

Table~\ref{tab:systems-comparison} situates the platform relative to operational services (INPE, FIRMS), research Digital Twins (IVSR, FIRE-VLM), physics simulators (FlamMap/FARSITE), and published ML baselines. The distinguishing combination is sub-minute LangGraph orchestration, three-class alert precision (82--92\%), and open-source deployment on a single EC2 instance---at the cost of state-specific calibration and limited GOES-16 production integration.

\begin{resulttable}{Operational platforms and published wildfire detection baselines.}{tab:systems-comparison}
\begin{minipage}[t]{0.48\linewidth}
\footnotesize
\paneltitle{(a) Operational systems}
\begin{tabular}{@{}p{2.0cm}p{1.1cm}ccp{1.6cm}@{}}
\toprule
\textbf{System} & \textbf{Lat.} & \textbf{Prec.} & \textbf{Ag.} & \textbf{Scope} \\
\midrule
INPE BD\-Queimadas & 3--6 h & ref. & No & Brazil \\
NASA FIRMS & 3--6 h & $\sim$80\% & No & Global \\
FIRE-VLM & sim. & --- & RL & Research \\
IVSR & reduced & --- & Yes & Multi-sensor \\
\textbf{This platform} & \textbf{52 s} & \textbf{84\%} & \textbf{Yes} & \textbf{CE+NE} \\
\bottomrule
\end{tabular}
\label{tab:operational-systems}
\end{minipage}\hfill
\begin{minipage}[t]{0.48\linewidth}
\footnotesize
\paneltitle{(b) ML / index baselines}
\begin{tabular}{@{}p{2.2cm}p{1.6cm}cc@{}}
\toprule
\textbf{System} & \textbf{Method} & \textbf{Prec.} & \textbf{Metric} \\
\midrule
Canadian FWI & Physical & --- & AUC 0.69 \\
Regional FWI+GA & Tailored & --- & AUC 0.79 \\
Decision Tree & ML/country & --- & AUC 0.86 \\
FC Autoencoder & Anomaly & --- & F1 0.74 \\
\textbf{This work} & GNN + Koopman & \textbf{82--92\%} & F1 0.54 \\
\bottomrule
\multicolumn{4}{@{}p{\linewidth}@{}}{\textit{Note:} Macro F1 from Table~\ref{tab:three-class-full}.} \\
\end{tabular}
\label{tab:sota}
\end{minipage}
\end{resulttable}
